\documentclass[../../../../SoftwareRequirementsSpecification.tex]{subfiles}
\begin{document}
% Richiede le macro definite in src/macros/useCaseMacros.tex
% (incluse nel preambolo di SoftwareRequirementsSpecification.tex)

\ucstart
\begin{tabularx}{\textwidth}{|l|X|}
  \hline
  \ucid{PT-02} \\ \hline
  \ucname{Registrazione nuovo Atleta} \\ \hline
  \ucactor{PT} \\ \hline
  \ucsecondaryactor{Atleta da Registrare} \\ \hline
  \ucdescription{Il PT crea l'account di un nuovo atleta e il
    sistema invita l'atleta ad accedere all'applicativo.} \\ \hline
  \ucprecond{Il PT deve essere loggato e deve
    avere la mail dell'atleta da registrare.} \\ \hline
  \ucpostcond{L'atleta ha un account nell'applicativo ed è stato
    invitato ad accedervi. Il suo profilo resta incompleto finché
    l'atleta non esegue il caso d'uso A-01.} \\ \hline
  \ucmainflow{
    \item Il PT atterra sulla pagina di creazione nuovo atleta.
    \item Il sistema mostra il form di inserimento dati dell'atleta
      (nome, cognome e mail).
    \item Il PT inserisce i dati.
    \item Il PT preme il pulsante "Crea atleta".
    \item Il sistema verifica che la mail sia valida e non sia già
      presente nel database.
    \item Il sistema crea l'account dell'atleta e lo associa al PT.
    \item Il sistema mostra una pagina di conferma creazione.
    \item Il sistema invia una mail all'atleta con l'invito ad
      accedere all'applicativo e le istruzioni per il primo accesso.
  } \\ \hline
  \ucaltflow{
    \textbf{5a. Email già presente} \newline
    - Il sistema mostra un messaggio di errore ("Atleta con questa
    email già presente"). \newline
    - Il flusso riprende dal punto 2. \newline
    \textbf{5b. Email non valida} \newline
    - Il sistema mostra un messaggio di errore ("Email non valida"). \newline
    - Il flusso riprende dal punto 2.
  } \\ \hline
  \ucnote{
    L'invito non contiene credenziali: l'atleta accede tramite il
    caso d'uso U-01 (Login), che gli invia un codice di accesso
    temporaneo, e completa poi il proprio profilo con A-01.
  } \\ \hline
\end{tabularx}
\end{document}
